\documentclass{article}
\usepackage{booktabs}
\usepackage[a4paper]{geometry}
\usepackage{fancyhdr}
\pagestyle{fancy}
\lhead{Komparative Kostenvorteile}
\rhead{März 2026}
\begin{document}
\section{Komparative Kostenvorteile}
Die \emph{absoluten Kostenvorteile} nach Adam Smith erklärten dass wenn ein Land $A$ in der Produktion aller Güterarten absolute produktiver ist als ein Land $B$, dass es dann keinen Grund für $A$ gäbe mit $B$ zu handeln. 
 
David Ricardo begründete mit der Theorie der \emph{komparativen Kostenvorteile} dass dies falsch ist. Sie dagt uns dass jedes Land sich auf die Produktion der Ware fokussieren sollte, welche es am effizientesten produzieren kann.
 
\subsection{Mathematische Begründung}
Angenommen beide Länder produzieren sowohl Ware $X$ als auch Ware $Y$. Dabei sind sie unterschiedlich effizient. Wird davon ausgegangen, dass mit dem Einsatz aller Arbeitskräfte diese Anzahl an Waren produziert werden kann. Die Einheit wird dabei außen vor gelassen. 
\begin{center}
 \begin{tabular}{l@{\hspace{2em}}cc}
  \multicolumn{3}{c}{Produktivität der Länder} \\
  \toprule
  & Land $A$ & Land $B$ \\
  \midrule
  Ware $X$ & 8 & 6 \\
  Ware $Y$ & 12 & 4 \\
  \bottomrule
 \end{tabular}
\end{center}
Wird sich nun Land $A$ angeschaut, so kann, der Theorie der absoluten Kostenvorteilen nach, mit voller Produktivität $8$ Einheiten von Ware $X$ oder $12$ Einheiten von Ware $Y$ oder etwas dazwischen Produzieren. 
\begin{center}
 \begin{tabular}{l@{\hspace{2em}}cccc}
  \multicolumn{5}{c}{Optionen der Produktion} \\
  \toprule
  & \multicolumn{2}{c}{Land $A$} & \multicolumn{2}{c}{Land $B$} \\
  \cmidrule{2-3} \cmidrule{4-5}
  & Ware $X$ & Ware $Y$ & Ware $X$ & Ware $Y$ \\
  \midrule
  Nur $X$ & 8 & 0 & 6 & 0\\
  Hälfte & 4 & 6 & 3 & 2\\  
  Nur $Y$ & 0 & 12 & 0 & 4\\
  \bottomrule
 \end{tabular}
\end{center}
Wird sich nun aber die Produktivität der Länder im Vergleich angesehen, so fällt auf dass Land $A$ produktiver in der Produktion von Ware $Y$ ist; für jede Ware $X$ könnten $1.5$ Ware $Y$ produziert werden. Land $B$ ist in der Produktion von Ware $Y$.
\begin{center}
 \begin{tabular}{ll@{\hspace{2em}}cc}
  \multicolumn{4}{c}{Wohlstand durch Handel} \\
  \toprule
  & & Land $A$ & Land $B$ \\
  \midrule
  Produktion & Ware $X$ & 0 & 6 \\
             & Ware $Y$ & 12 & 0 \\
  \midrule
  Handel     & Ware $X$ & 4 & -4 \\
             & Ware $Y$ & -4 & 4 \\
  \midrule
  Wohlstand  & Ware $X$ & 4 & 2 \\
             & Ware $Y$ & 8 & 4 \\
  \bottomrule
 \end{tabular}
\end{center}
Während vorher ausgerechnet wurde dass Land $A$, wenn es 4 Einheiten an Ware $X$ wollte, konnte höchstens 6 Einheiten an Ware $Y$ produzieren. Nun kann es bei gleicher Anzahl $X$ zwei weitere Einheiten $Y$ produzieren. Gleiches kann für Land $B$ gesagt werden.
 
Dass sich beide Länder auf ihre produktivste Ware fokussierten und handelten war zum vorteil beider Länder.
\end{document}